% !PRINTED 84
\weberpage{84}
Ist diese Bedingung erfüllt, so ist gewiss $g_1, g_2, \dots g_\nu$ die Basis einer Gruppe mit den Invarianten $i_1, i_2, \dots i_\nu$. Wenn aber auch noch die in IV. aufgestellte Forderung erfüllt sein soll, dann folgt noch weiter:

\gesperrt{3. Ist $q_h$ eine in $m$ aufgehende Primzahl, der der Index $\beta_h$ entspricht, so dürfen die $\nu$ Grössen}
\[
\gamma_{1,h}, \gamma_{2,h}, \dots \gamma_{\nu,h},
\]
\gesperrt{wenn $q_h$ mehrfach in $m$ aufgeht, nicht alle durch $q_h$, und wenn $q_h$ nur einmal in $m$ aufgeht, nicht alle durch $q_h-1$ theilbar sein.}

Die Bedingungen 2., 3. sind dann also die nothwendigen und hinreichenden, und es kommt jetzt nur noch darauf an, sie in eine Form zu bringen, dass die Möglichkeit ihrer Erfüllung beurtheilt werden kann.

Wir bezeichnen mit $\delta_{k,h}$ \gesperrt{den grössten gemeinschaftlichen Theiler von $i_k$ und $c_h$ und setzen}
\weq{11}{i_k=\delta_{k,h}\,i_{k,h}, \quad c_h=\delta_{k,h}\,c_{k,h}, \quad \begin{aligned}k&=1,2,\dots\nu,\\ h&=1,2,\dots\mu,\end{aligned}}
so dass $i_{k,h}$ und $c_{k,h}$ \gesperrt{relativ prim} sind. Es ist dann zu bemerken, dass die Grössen
\[
\delta_{k,h}, \quad i_{k,h}, \quad c_{k,h}
\]
durch $m$ und die Invarianten $i_k$ völlig bestimmt sind. Nach der Definition der Invarianten müssen die $\delta_{k,h}$ und $i_{k,h}$ \gesperrt{Primzahlpotenzen} sein.

Aus den Voraussetzungen, die wir in a) bis d) über die Zahl $m$ gemacht haben, ergiebt sich eine Folgerung für die $c_{k,h}$, die wir hervorheben müssen.

\gesperrt{4. Ist $q_h^\varkappa$ einer der Factoren von $m$, der dem Index $\beta_h$ entspricht, ist also $c_h=q_h^{\varkappa-1}(q_h-1)$, oder wenn $q=2$ und $\varkappa>2$ ist, $c_h=2^{\varkappa-2}$, und wenn $\varkappa=2$ ist, $c_h=2$, so können, wenn $\varkappa>1$ ist, die Zahlen}
\weq{12}{c_{1,h}, c_{2,h}, \dots c_{\nu,h}}
\gesperrt{nicht alle durch $q_h$, und wenn $\varkappa=1$ ist, nicht alle durch $q-1$ theilbar sein.}
% [sic: Weber switches from q_h to q in the exceptional 2-adic clause and in the final q-1 clause; q_h is intended in both places.]

Wenn nämlich die Grössen (12) alle durch $q_h$ theilbar sind, so sind nach (11) die Zahlen
\weq{13}{\delta_{1,h}.\ \delta_{2,h}, \dots \delta_{\nu,h}}
% [sic: Weber prints a full stop, rather than a comma, after the first term.]
%CONT>>
